\documentclass[../book.tex]{subfiles}
\begin{document}

    \section{一元函数微分学的概念}

    \subsection{导数}

    \begin{enumerate}
        \item 设\( y = f(x) \)定义在区间\( I \)上，让自变量在\( x = x_0 \)处加一个增量\( \Delta x \)（可正可负），其中\( x_0 \in I, x_0 + \Delta x \in I \)，则可得函数的增量\( \Delta y = f(x_0 + \Delta x) - f(x_0) \)。若函数增量\( \Delta y \)与自变量增量\( \Delta x \)的比值在\( \Delta x \to 0 \)时的极限存在，即\( \lim\limits_{\Delta x \to 0} \displaystyle\frac{\Delta y}{\Delta x} \)存在，则称函数\( y = f(x) \)在点\( x_0 \)处可导，并称这个{\color{red}{极限}}为\( y = f(x) \)在点\( x_0 \)处的{\color{red}{导数（变化率）}}，记作\( f'(x_0) \)，即
        \[
            \frac{dy}{dx}\bigg|_{x=x_0} = f'(x_0) = \lim\limits_{\Delta x \to 0} \frac{\Delta y}{\Delta x} = \lim\limits_{\Delta x \to 0} \frac{f(x_0 + \Delta x) - f(x_0)}{\Delta x}.
            \tag{*}
        \]
        广义化
        \[
            \lim_{\text{狗} \to 0} \frac{f(x_0 + \text{狗}) - f(x_0)}{\text{狗}}.
            \tag{**}
        \]
        令 $x_0 + \Delta x = x$，从而得到
        \[
            \text{函数式:\quad} f'(x_0) = \lim_{x \to x_0} \frac{f(x) - f(x_0)}{x - x_0},
            \tag{***}
        \]
        \item 下面这三种提法是等价的:
        \begin{itemize}
            \item $y = f(x)$在点$x_0$处可导；
            \item $y = f(x)$在点$x_0$处导数存在；
            \item $f'(x_0) = A$（$A$为有限数）。
        \end{itemize}
        \item $\color{red}{\bigstar}$ 函数在一点可导的充要条件: $f'(x_0)$存在 $\Leftrightarrow$ 其左导数$f'_-(x_0)$与右导数$f'_+(x_0)$均存在且相等
        \item 函数在一点可导的必要条件: \textbf{若$f(x)$在一点可导 $\Rightarrow$ $f(x)$在该点连续}，反之未必
        \item 导数的性质
        \begin{itemize}
            \item $\color{red}{\bigstar}$求导一次，奇偶性互换（导数定义）
            \item 若$f(x)$是可导的周期为$T$的周期函数，则$f'(x)$也是以$T$为周期的周期函数（导数定义）
            \item 若$f'(x) > 0$，$f(x)$单调递增；若$f'(x) < 0$，则$f(x)$单调递减
        \end{itemize}
    \end{enumerate}

    \subsection{导数的几何意义}

    \begin{enumerate}
        \item 函数 $y=f(x)$ 在点 $x_0$ 处的导数值 $f'(x_0)$，就是曲线 $y=f(x)$ 在点 $(x_0,\,y_0)$处切线的斜率 $k$，即 $k = f'(x_0)$，因此，曲线 $y=f(x)$ 在点 $(x_0,\,y_0)$ 处的切线方程为
        \[
            y - y_0 = f'(x_0)\,(x - x_0).
        \]

        该点处的法线方程为
        \[
            y - y_0 = -\frac{1}{f'(x_0)}\,(x - x_0) \; (f'(x_0) \neq 0)
        \]
        \item 切线存在不代表导数存在，但{\color{red}{导数存在切线一定存在}}
        \begin{itemize}
            \item $f'_+(x_0) \neq f'_-(x_0)$，出现角点（尖点），则$f(x)$在$x_0$处不可导，没有切线。\emph{e.g.} $f(x) = |x|$在$x = 0$处
            \item $f(x)$在点$x_0$的导数是无穷导数时，在该点有切线但无导数。\emph{e.g.} $f(x) = x^\frac{1}{3}$在$x = 0$处
        \end{itemize}
    \end{enumerate}

    \subsection{微分的概念}

    \begin{enumerate}
        \item 设函数$y = f(x)$在点$x_0$的某邻域内有定义，且$x_0 + \Delta x$在该邻域内，对于增量函数

        \noindent
        \begin{minipage}{0.54\textwidth}
            \[
                \color{red}{\Delta y = f(x_0 + \Delta x) - f(x_0)}
            \]

            若存在与$\Delta x$无关的常数$A$，使得
            \[
                \color{red}{
                    \Delta y = A\,\Delta x + o(\Delta x)
                    = dy + o(\Delta x)
                }
            \]
        \end{minipage}
        \hfill
        \begin{minipage}{0.44\textwidth}
            \centering
            \includegraphics[width=\linewidth]{\subfix{../../../assets/math/advanced-mathematics/pic3-15.png}}
        \end{minipage}

        其中$o(\Delta x)$是在$\Delta x \to 0$时比$\Delta x$更高阶的无穷小，则称$f(x)$在点$x_0$处\textbf{可微}，并把增量的主要部分$A\Delta x$称为线性主部，也叫做$f(x)$在点$x_0$处的微分，记
        \[
            \color{red}{
                \left. dy \right|_{x=x_0} = A\Delta x = f'(x_0)\,dx
            }
        \]
        \item 可微的判别
        \begin{itemize}
            \item 一元函数可微 $\Leftrightarrow$ 可导
            \item 做极限
            \begin{circlenum}
                \item 写增量 $\Delta y = f(x_0 + \Delta x) - f(x_0)$
                \item 写线性增量 $A\Delta x = f'(x_0)\Delta x$
                \item 做极限
                \[
                    \lim_{\Delta x \to 0}
                    \frac{\Delta y - A\Delta x}{\Delta x}
                    \Leftrightarrow
                    \Delta y = A\Delta x + o(\Delta x)
                \]
                \item 若该极限等于 $0$，则 $y=f(x)$ 在点 $x_0$ 处可微，否则不可微
            \end{circlenum}
        \end{itemize}
        \item 可微的含义：用形式简单的“线性增量 $A\Delta x$”去代替形式复杂的“增量 $\Delta y$”，且其误差“$\Delta y - A\Delta x$”是 $o(\Delta x)$，即用“简单的量”代替了“复杂的量”，且产生的误差又可以忽略不计
        \item 可微的几何意义：若$f(x)$在点$x_0$处可微，则在点$(x_0, y_0)$附近可以用切线段近似代替曲线段
    \end{enumerate}

    \subsection{定理}

    \begin{enumerate}
        \item \textbf{达布定理：} 若函数 $f(x)$ 在 $[a,b]$ 上可导，且 $f'_+(a)\neq f'_-(b)$，则对于任意介于 $f'_+(a)$ 与 $f'_-(b)$ 之间的数 $u$，存在 $\xi\in(a,b)$，使得
        \[
            f'(\xi)=u.
        \]
    \end{enumerate}

    \subsection{结论}

    \begin{enumerate}
        \item $\color{red}{\bigstar}\color{red}{\bigstar}$两个函数的乘积中，如果一个因子在某点连续但不可导，而乘积函数在该点可导，那么另一个因子在该点必定为零
        \item 若函数 $\varphi(x)$ 在 $x=x_0$ 处连续，则函数$f(x)=|x-x_0|\,\varphi(x)$在 $x=x_0$ 处可导的充要条件是 $\varphi(x_0)=0$
        \item {\color{red}{零点可以{\blueuline{压制}}不可导性}}。不可导意味着函数在那一点有{\blueuline{尖角}}（斜率突变）。但乘上一个在该点为零的因子后，乘积函数在附近会被{\blueuline{拉向零}}，尖角就被{\blueuline{压平}}了。用一个递进的例子看：
        \begin{itemize}
            \item $|x|$：在 $0$ 处有尖角，不可导
            \item $x \cdot |x|$：乘了 $x$（在 $0$ 处为零），尖角被压平了，可导
            \item $x^2 \cdot |x|$：乘了 $x^2$（在 $0$ 处为零且{\blueuline{压得更狠}}），不仅可导，导数还更光滑
        \end{itemize}
        \item 函数 $f(x)$ 与 $|f(x)|$ 的连续性与可导性关系总结：
        \begin{itemize}
            \item 若 $f(x)$ 在 $x_0$ 处连续，则 $|f(x)|$ 在 $x_0$ 处连续； 反之不成立。
            \item 若 $f(x)$ 在 $x_0$ 处可导，则：
            \begin{enumerate}[label=(\alph*)]
                \item 若 $f(x_0) \neq 0$，则 $|f(x)|$ 在 $x_0$ 处{\color{red}{可导}}，且
                \[
                    \left.(|f(x)|)'\right|_{x=x_0}
                    =
                    \begin{cases}
                        f'(x_0), & f(x_0) > 0, \\[0.3em]
                        -f'(x_0), & f(x_0) < 0.
                    \end{cases}
                \]

                \item 若 $f(x_0) = 0$，则：
                \[
                    \begin{cases}
                        f'(x_0) = 0
                        \;\Rightarrow\;
                        |f(x)| \text{ 在 } x_0 \text{ 处可导，且 }
                        \left.(|f(x)|)'\right\Big|_{x=x_0} = 0, \\[0.5em]
                        f'(x_0) \neq 0
                        \;\Rightarrow\;
                        |f(x)| \text{ 在 } x_0 \text{ 处{\color{red}{不可导}}}.
                    \end{cases}
                \]
                \item 总结不可导即:
                \[
                    \left.
                    \begin{aligned}
                        & f(x) \text{ 在 } x_0 \text{ 处可导} \\
                        & f(x_0) = 0 \\
                        & f'(x_0) \neq 0
                    \end{aligned}
                    \right\}
                    \;\Longrightarrow\;
                    |f(x)| \text{ 在 } x_0 \text{ 处 {\color{red}{不可导}}}
                \]
            \end{enumerate}
        \end{itemize}
        \item $f^{(n)}(x_0)$存在 $\Rightarrow f^{(n - 1)}(x)$在$x_0$附近有定义且在$x_0$处连续
    \end{enumerate}

    \subsection{理解}

    \begin{enumerate}
        \item 函数 $f(x)$ 存在与 $f'(x)$ 存在的区别
        \begin{enumerate}
            \item 函数 $f(x)$ 在 $x=x_0$ 处的极限存在，不能推出 $f(x)$ 在 $x=x_0$ 处连续。例如，$x_0$ 可以是 $f(x)$ 的可去间断点
            \item 若 $f(x)$ 在 $x_0$ 处可导，且$\lim\limits_{x\to x_0} f'(x)=a$，则 $f'(x)$ 在 $x_0$ 处连续。即
            \[
                \lim_{x\to x_0} f'(x)=a,\quad f'(x_0)=b
                \;\Longrightarrow\;
                \lim_{x\to x_0} f'(x)=f'(x_0)=a=b.
            \]
            \item $f(x)$存在不能得出$f(x)$有介值性。因为“存在”不等于“连续”，而只有连续函数才具有介值性
            \item $f'(x)$存在，可得$f'(x)$有介值性
            \item $f'(x)$存在且$\neq 0$，则$f'(x)$必保号（恒正或恒负）
            \begin{analysisbox}[证]
                用反证法。
                若存在 $a,b$ 使
                \[
                    f'(a)<0,\quad f'(b)>0,
                \]
                则由达布定理，存在 $\xi\in(a,b)$，使得
                \[
                    f'(\xi)=0,
                \]
                这与 $f'(x)\neq0$ 的假设矛盾
            \end{analysisbox}
            \item 若 $f'(x)$ 在 $[a,b]$ 上存在，则 $f'(x)$ 在该区间内\textbf{不存在第一类间断点}
        \end{enumerate}
        \item 对于连续或可导函数，只要$f(x) > 0$（或$f(x) < 0$），无论$f(x_0)$与$0$的距离有多小，它旁边相依相偎的$f(x)$一定$> 0$（或$< 0$）

        \item 求导数时，当函数不具备{\color{red}{导数存在}}的条件时，往往只能用{\color{red}{导数定义}}求
        \begin{analysisbox}[Example]
            设$f(x) = (x-a)\cdot\varphi(x)$，其中 $\varphi(x)$ 连续，求 $f'(a)$。

            \[
                \begin{aligned}
                    f'(a)
                    &= \lim_{x\to a}\frac{f(x)-f(a)}{x-a} \\
                    &= \lim_{x\to a}\frac{(x-a)\cdot\varphi(x)-0}{x-a} \\
                    &= \lim_{x\to a}\cdot\varphi(x) \\
                    &= \varphi(a)
                \end{aligned}
            \]
        \end{analysisbox}

        \item \begin{theorem}
                  证明：
                  (1)\ 若 $F(x)$ 在 $[x_0,\,x_0+\delta)\ (\delta>0)$ 上连续，在 $(x_0,\,x_0+\delta)$ 内可导，当 $\lim\limits_{x\to x_0^+} F'(x) = A$ 存在时，有 $F'_+(x_0) = A$

                  \medskip

                  (2)\ 若 $F(x)$ 在 $(x_0-\delta,\,x_0]\ (\delta>0)$ 上连续，在 $(x_0-\delta,\,x_0)$ 内可导，当 $\lim\limits_{x\to x_0^-} F'(x) = A$ 存在时，有 $F'_-(x_0) = A$
        \end{theorem}
        \begin{analysisbox}[证明]
            \[
                \begin{aligned}
                    \text{(1)\quad} F'_+(x_0)
                    &= \lim_{x \to x_0^+} \frac{F(x) - F(x_0)}{x - x_0}
                    \overset{\text{洛必达法则}}{=} \lim_{x \to x_0^+} \frac{F'(x)}{1} = A,\\[0.5em]
                    \text{(2)\quad} F'_-(x_0)
                    &= \lim_{x \to x_0^-} \frac{F(x) - F(x_0)}{x - x_0}
                    \overset{\text{洛必达法则}}{=} \lim_{x \to x_0^-} \frac{F'(x)}{1} = A.
                \end{aligned}
            \]
        \end{analysisbox}

        \item 设 $f(x)$ 在 $x=0$ 的某邻域内有定义，且 $|f(x)| \le 1-\cos x$，证明 $f(x)$ 在 $x=0$ 处连续、可导且 $f'(0)=0$
        \begin{analysisbox}[经典]
            \textbf{证明：} 分三个层次完成。

            \medskip
            \textcolor{analysisTitleColor}{\textbf{（1）求函数极限}}

            由
            \[
                0 \le |f(x)| \le 1-\cos x,
                \qquad
                \lim_{x\to 0}(1-\cos x)=0,
            \]
            根据夹逼准则，有
            \[
                \lim_{x\to 0}|f(x)| = 0
                \quad\Rightarrow\quad
                \lim_{x\to 0}f(x) = 0.
            \]

            \medskip
            \textcolor{analysisTitleColor}{\textbf{（2）求函数值，验证连续性}}

            在原不等式中令 $x=0$，得
            \[
                |f(0)| \le 1-\cos 0 = 0,
            \]
            故 $f(0)=0$。

            因此
            \[
                \lim_{x\to 0}f(x) = f(0),
            \]
            即 $f(x)$ 在 $x=0$ 处连续，并且有
            \[
                |f(x)-f(0)| = |f(x)| \le 1-\cos x.
            \]

            \medskip
            \textcolor{analysisTitleColor}{\textbf{（3）求导数（按定义）}}

            由上述不等式可得
            \[
                0 \le
                \left|
                \frac{f(x)-f(0)}{x-0}
                \right|
                \le
                \frac{1-\cos x}{|x|}.
            \]

            又因为
            \[
                \lim_{x\to 0}\frac{1-\cos x}{|x|}
                =
                \lim_{x\to 0}\frac{\frac{1}{2}x^2}{|x|}
                =0,
            \]
            再次由夹逼准则，
            \[
                \lim_{x\to 0}
                \left|
                \frac{f(x)-f(0)}{x-0}
                \right|
                =0,
            \]
            即
            \[
                \lim_{x\to 0}
                \frac{f(x)-f(0)}{x-0}
                =0,
            \]
            从而 $f'(0)=0$。

            \medskip
            \textbf{综上，}
            $f(x)$ 在 $x=0$ 处连续、可导，且 $f'(0)=0$

        \end{analysisbox}
    \end{enumerate}
\end{document}
